\section{Conclusions}\label{sec:conclusions}
A production-adjusted Ridge model gave a more accurate and more traceable baseline for pelleting electricity than the operational references evaluated here. In the frozen chronological test it reduced MAE by 14.42\% relative to line-specific SEC while providing a reproducible additive decomposition of every estimate. The improvement concerns the estimation of expected electricity under the recorded production conditions; it is not a reduction in plant energy use or evidence of achievable savings.

Under the tested 16-record, same-line, within-week representation, completed operating histories did not improve on current-record covariates, and neither TFT, N-HiTS nor the tested TimesFM~2.5 and Chronos-2 configurations beat the current-input control. Production adjustment helped on all three lines, most consistently on Pellet~2 and Pellet~3; Pellet~1 improved on SEC, but its systematic underestimation calls for line-specific recalibration or additional predictors before prospective use.

The evidence supports retrospective assessment at the closure of recorded operations, where the baseline provides an expected value, an empirically calibrated residual interval and a traceable deviation for identifying candidates for technical review. Prospective supervisory use requires validation on a new consecutive production period with prespecified line-level criteria and verified record identity, measurement boundaries, predictor availability and calibration rules. Until then, automatic control, modification of process set-points and operational digital-twin integration remain outside the demonstrated scope of the study.
